\section{牛顿力学的基本定律 II \& 动量}

\subsection{惯性系与非惯性系}

牛顿运动定律成立的参考系是惯性参考系。牛顿运动定律不成立的参考系是非惯性参考系。

相对于已知惯性系作匀速直线运动的参考系也是惯性系，但是作加速运动的参考系则是非惯性系。

\subsection{非惯性系中的惯性力}

惯性力是为了让牛顿运动定律在非惯性系下成立而引入的力。

\subsubsection{匀加速非惯性系中的惯性力}

设非惯性系 $S'$ 相对于惯性系 $S$ 以加速度 $\vect{a}$ 运动，那么惯性力为：
$$
\vect{F} = -m \vect{a}
$$

\subsubsection{匀速转动非惯性系中的惯性力}

\begin{enumerate}
    \item 物体在非惯性系中静止时，有惯性离心力：
    $$
    \vect{F}_r = m \vect{r} \omega^2 = m (\vect{\omega} \times \vect{r}) \times \vect{\omega} = m \vect{v}_{\theta} \times \vect{\omega}
    $$
    \item 物体在非惯性系中运动时，有科里奥利力。设物体 $m$ 在 $S'$ 中沿径向匀速运动，在惯性系中位矢为 $\vect{r}(t)$，径向速度为 $\vect{v}_r(t)$，横向速度为 $\vect{v}_\theta(t) = \vect{\omega} \times \vect{r}$。那么在 $S$ 中：
    $$
    \begin{aligned}
        \vect{v} & = \vect{v}_r + \vect{v}_{\theta} \\
        \vect{v}_{\theta} & = \vect{\omega} \times \vect{r} \\
        \vect{a} & = \dv{\vect{v}}{t} = \dv{\vect{v}_r}{t} + \dv{\vect{v}_{\theta}}{t} \\
        & = \dv{v_r}{t} \vect{e}_r + v_r \dv{\vect{e}_r}{t} + \dv{\vect{\omega}}{t} \times \vect{r} + \vect{\omega} \times \dv{\vect{r}}{t} \\
        & = \vect{a}' + \vect{\omega} \times \vect{v}_r + \vect{\omega} \times \ab(\vect{v}_{r} + \vect{v}_{\theta}) \\
        & = \vect{a}' + 2 \vect{\omega} \times \vect{v}_r + \vect{\omega} \times \vect{v}_{\theta}
    \end{aligned}
    $$
    其中第一项是 $S'$ 中的加速度，第三项是惯性离心力，第二项是科里奥利力：
    $$
    \vect{F}_C = 2 m \vect{v}_r \times \vect{\omega}
    $$
\end{enumerate}

\subsubsection{一般转动非惯性系中的惯性力}

设物体 $m$ 在 $S'$ 中速度为 $\vect{v}'$，加速度为 $\vect{a}'$。那么在 $S$ 中：
$$
\vect{a}_S = \vect{a}' + \dv{\vect{\omega}}{t} \times \vect{r} - (\vect{\omega} \times \vect{r}) \times \vect{\omega} - 2 \vect{v}' \times \vect{\omega}
$$
得到惯性力包括：
\begin{itemize}
    \item 惯性离心力：
    $$
    \vect{F}_r = m (\vect{\omega} \times \vect{r}) \times \vect{\omega}
    $$
    \item 科里奥利力：
    $$
    \vect{F}_C = 2 m \vect{v}' \times \vect{\omega}
    $$
    \item 欧拉力：
    $$
    \vect{F}_E = -m \dv{\omega}{t} \times \vect{r}
    $$
\end{itemize}

\section{动量}

\begin{definition}[动量和冲量]
    质点的\textbf{动量} $\vect{p}$ 定义为 $m \vect{v}$。动量的变化量称为\textbf{冲量}。
\end{definition}

根据牛顿运动定律，可以直接得到质点的动量变化定理：

\begin{theorem}[动量变化定理]
    $$
    \dd \vect{p} = \vect{F} \dd t
    $$
    或者积分形式：
    $$
    \vect{I} = \int_{t_1}^{t_2} \vect{F} \dd t = \vect{p}_2 - \vect{p}_1
    $$
\end{theorem}

动量变化定理在微观、高速的情况下仍然成立
